\PilotPage{18}
\begin{equation}
    m\ddot{u}(t)+k u(t)=0.
    \tag{4}
\end{equation}
Eq.~(4) is a linear homogeneous ordinary differential equation (ODE).

\subsubsection{Solution of Linear ODEs}
The solution of linear homogeneous ODEs is sought in the form $e^{st}$.
\begin{equation}
    % SOURCE CORRECTION: The dependent variable is consistently $u(t)$ in this derivation.
    u(t)=e^{st}.
    \tag{5}
\end{equation}
Upon differentiation, Eq.~(5) gives
\begin{equation}
    \begin{aligned}
        \dot{u}(t)&=s e^{st}=s u(t),\\
        \ddot{u}(t)&=s^2 e^{st}=s^2 u(t).
    \end{aligned}
    \tag{6}
\end{equation}
Substitution of Eqs.~(5), (6) into Eq.~(4) yields
\begin{equation}
    m s^2 u(t)+k u(t)=0.
    \tag{7}
\end{equation}
After factoring out $u(t)$, Eq.~(7) becomes
\begin{equation}
    (m s^2+k)u(t)=0.
    \tag{8}
\end{equation}
A possible solution of Eq.~(8) is $u(t)=0$. However, this solution is trivial and of no interest to us. We are interested in nontrivial solutions $u(t)\ne0$. To obtain a nontrivial solution $u(t)\ne0$, one needs to find the values of $s$ that make zero the factor multiplying $u(t)$ in Eq.~(8), i.e., to have $(m s^2+k)=0$.

\subsubsection{Characteristic Equation}
% SOURCE CORRECTION: The referenced subsection in this chapter is 1.1.4.
It was shown in Section 1.1.4 that a nontrivial solution $u(t)\ne0$ of Eq.~(8) can only be obtained if the factor multiplying $u(t)$ is zero, i.e.,
\begin{equation}
    m s^2+k=0.
    \tag{9}
\end{equation}
Divide Eq.~(9) by $m$ to get what is called the characteristic equation, i.e.,
\begin{equation}
    s^2+\frac{k}{m}=0 \qquad \text{(characteristic equation)}.
    \tag{10}
\end{equation}
Solution of the characteristic equation will give us the values of $s$ that permit nontrivial solution $u(t)\ne0$ of Eq.~(8). Upon solution, Eq.~(10) gives
\begin{equation}
    s^2=-\frac{k}{m}.
    \tag{11}
\end{equation}
Denote
\begin{equation}
    \omega_n^2=\frac{k}{m}, \qquad \omega_n=\sqrt{\frac{k}{m}} \quad \text{(natural frequency)}.
    \tag{12}
\end{equation}
The natural frequency $\omega_n$ is expressed in rad/sec. Substitution of Eq.~(12) into Eq.~(11) yields
